% !TEX root = ../main.tex

\chapter{TINJAUAN PUSTAKA}

\section{Penelitian Terdahulu}
Dibandingkan dengan lima penelitian pada Tabel~\ref{tab:penelitian_terdahulu},
studi ini berfokus pada pengembangan sistem \textit{end-to-end}. Sistem yang diusulkan
mengintegrasikan deteksi, pelacakan, dan analitik pelanggaran helm secara
\textit{real-time} dengan model YOLOv8 dan arsitektur \textit{big data streaming}.
Pendekatan ini berbeda dari studi terdahulu yang umumnya hanya membahas deteksi
atau pelacakan. Luaran penelitian ini adalah alur analitik yang menghubungkan hasil deteksi,
pelacakan, dan penyimpanan data secara otomatis.

\vspace{0.5\baselineskip}

{\footnotesize
  \setlength{\tabcolsep}{4pt}
  \begin{longtable}{@{}>{\RaggedRight\arraybackslash}p{0.23\textwidth} >{\RaggedRight\arraybackslash}p{0.25\textwidth} >{\RaggedRight\arraybackslash}p{0.14\textwidth} >{\RaggedRight\arraybackslash}p{0.30\textwidth}@{}}
    \caption{Ringkasan Penelitian Terdahulu}
    \label{tab:penelitian_terdahulu} \\
    \toprule
    \textbf{Peneliti \& Judul} & \textbf{Metodologi} & \textbf{Data} & \textbf{Hasil Utama} \\
    \midrule
    \endfirsthead

    \caption[]{Ringkasan Penelitian Terdahulu} \\
    \toprule
    \textbf{Peneliti \& Judul} & \textbf{Metodologi} & \textbf{Data} & \textbf{Hasil Utama} \\
    \midrule
    \endhead

    \midrule
    \endfoot

    \bottomrule
    \endlastfoot

    \textbf{Figo Agil Alunjati \cite{alunjati2025intersection}} \newline
    \textit{Architecture of Four-Arm Intersection Performance Evaluation System Based on Turn Movement Count and Omnidirectional Camera.}(2025) &
    Sistem evaluasi simpang \textit{real-time} berbasis kamera omnidirectional dan YOLO (v5, v8, v11) untuk \textit{turn movement count}. &
    CCTV ATCS Sedayu (2 hari). &
    YOLOv11 (batch 32): \textbf{mAP 93\%, F1 91\%, FPS 53}; \textit{occlusion} berkurang, tetapi sensitif terhadap cahaya ekstrem. \\

    \textbf{Liu et al. \cite{liu2025vehicle}} \newline
    \textit{Vehicle Flow Detection dengan YOLOv8n + ByteTrack.}(2025) &
    Integrasi MSN-YOLOv8n (SEAM, NWD Loss) dan ByteTrack--ReID untuk deteksi serta pelacakan kendaraan kecil. &
    UA-DETRAC. &
    \textbf{mAP 62.8\%}, \textbf{MOTA 72.16\%}; stabil untuk objek kecil, tetapi baru diuji pada satu dataset dan belum \textit{streaming real-time}. \\

    \textbf{Aboah et al. \cite{aboah2023helmet}} \newline
    \textit{Real-Time Multi-Class Helmet Violation Detection Using Few-Shot Data Sampling Technique and YOLOv8.}(2023) &
    YOLOv8 dengan \textit{few-shot data sampling} untuk deteksi pelanggaran helm multikelas. &
    AI City 2023 Track 5. &
    YOLOv8+TTA: mAP 0.5861, 95 FPS, peringkat 7/192; efisien untuk implementasi \textit{real-time} \textit{smart city}. \\

    \textbf{Do et al. \cite{do2025ramots}} \newline
    \textit{RAMOTS: Real-Time Aerial Multi-Object Tracking.}(2025) &
    Pelacakan multiobjek UAV dengan YOLOv8/v10 + BYTETrack pada arsitektur Kafka--Spark. &
    Video UAV. &
    \textbf{HOTA 48.14}, \textbf{28 FPS}; efektif untuk udara, namun menurun di lingkungan darat dan masih bergantung satu GPU. \\

\end{longtable}}

\section{Landasan Teori}

\subsection{Big Data}

\textit{Big Data} merupakan kumpulan data berukuran besar, tumbuh cepat, dan
beragam. Karakteristik ini membuat pengolahan tradisional
kurang efisien. Karakteristik utama \textit{Big Data} dikenal sebagai "5V"
\cite{marz2015big}:

\begin{enumerate}[a.]
  \item \textit{Volume}: Data berukuran besar, misalnya aliran video 24/7 dari kamera CCTV dan data pelanggaran yang terus bertambah pada \textit{data warehouse}. Kondisi ini membutuhkan infrastruktur penyimpanan dan pemrosesan yang skalabel.

  \item \textit{Velocity}: Data perlu diproses secara \textit{real-time} dengan target latensi $\leq$ 2 detik, sehingga arsitektur sistem harus efisien dan stabil.

  \item \textit{Variety}: Data tersedia dalam berbagai format:
  \begin{enumerate}[1).]
    \item Tidak terstruktur: video, gambar, dan media sosial.
    \item Semi-terstruktur: log sistem dari Apache Kafka.
    \item Terstruktur: data pelanggaran dalam skema relasional.
  \end{enumerate}

\item \textit{Veracity}: Keandalan dan akurasi data harus dijaga melalui proses validasi untuk menekan \textit{false positive} dan \textit{false negative}.

\item \textit{Value}: Data perlu menghasilkan nilai praktis, misalnya dalam bentuk \textit{dashboard} analitik untuk mendukung pengambilan keputusan.
\end{enumerate}

Penelitian ini menghadapi tantangan utama pada aspek \textit{volume} dan
\textit{velocity}. Kondisi tersebut menjadi alasan penggunaan
\textit{stream processing} dibandingkan \textit{batch processing}. Salah satu
pendekatan yang relevan adalah \textit{Kappa Architecture}, yang memungkinkan
pemrosesan historis dan \textit{real-time} melalui satu alur data
\cite{marz2015big}.

\subsection{Arsitektur Streaming: \textit{Kappa Architecture}}

Arsitektur Kappa adalah pendekatan pemrosesan \textit{Big Data} yang berfokus
pada aliran data \textit{real-time} \cite{kreps2014iheartlogs}. Pendekatan
ini menyederhanakan arsitektur Lambda dengan menghilangkan pemisahan antara
\textit{batch layer} dan \textit{speed layer}.  Data bersifat temporal dan dapat dipandang sebagai
aliran peristiwa (\textit{event stream}). Karena itu, paradigma
\textit{stream-first} memproses data historis dan data baru melalui jalur yang
sama. Pada arsitektur ini, repositori log seperti Apache~Kafka berperan
sebagai \textit{single source of truth} dan tetap memungkinkan proses
\textit{replay} saat diperlukan. Dengan satu alur pemrosesan, arsitektur
Kappa digunakan untuk menekan latensi dan menjaga konsistensi implementasi
\cite{kreps2014questioning}. Dengan rancangan ini, pengembangan, pengujian, dan
pemeliharaan sistem analitik menjadi lebih sederhana.

Arsitektur Kappa terdiri atas tiga komponen utama:
\begin{enumerate}[a.]
\item \textit{Ingestion Layer}: menerima data dari berbagai sumber dan memasukkannya ke antrean data menggunakan \textit{message broker} seperti Apache~Kafka.
\item \textit{Stream Processing Layer}: melakukan transformasi, agregasi, dan analisis data secara berkelanjutan menggunakan \textit{stream processing framework}.
\item \textit{Serving Layer}: menyimpan hasil olahan pada sistem penyimpanan yang mendukung \textit{low-latency read} untuk kebutuhan pengguna maupun sistem lain.
\end{enumerate}

Dengan satu \textit{pipeline}, Arsitektur Kappa cocok untuk aplikasi berlatensi rendah seperti \textit{IoT analytics} dan pemantauan lalu lintas. Pendekatan ini mengurangi kebutuhan sinkronisasi antarlapisan, mempercepat waktu respons, dan menyederhanakan pemeliharaan sistem.

\begin{figure}[H]
\centering
\includegraphics[width=1\textwidth]{gambar/Kappa-architecture-data-flow.png}
\caption{Alur Data Arsitektur Kappa \cite{kreps2014questioning}}
\label{fig:lambda_vs_kappa}
\end{figure}

\subsubsection{Apache Kafka}
Apache Kafka adalah platform \textit{streaming} terdistribusi yang dirancang
untuk menangani aliran data berkecepatan tinggi dengan latensi rendah dan
keandalan yang baik. Dalam arsitektur data \textit{real-time}, Kafka berperan
sebagai \textit{message broker} yang menghubungkan produsen dan konsumen data
secara asinkron, termasuk pada sistem berskala besar berbasis
\textit{microservices} \cite{kreps2014questioning}. Penerapan Kafka pada
pemrosesan video \textit{real-time} juga telah dilaporkan dalam penelitian
\cite{ayae2017video}, yang menunjukkan bahwa Kafka mampu menjaga aliran data tetap berjalan.

Kafka menyimpan data sebagai rekaman berurutan di dalam
\textit{topic} yang dibagi ke beberapa \textit{partition}. Mekanisme ini
memungkinkan proses tulis-baca berlangsung paralel, sehingga \textit{throughput}
sistem meningkat tanpa mengorbankan latensi. Setiap data memiliki
\textit{offset} sebagai penanda posisi, sehingga status konsumsi data dapat
dilacak secara tepat dan diputar ulang (\textit{replay}) saat dibutuhkan.

Kafka juga menyediakan reliabilitas melalui replikasi
\textit{partition} antar-\textit{broker} untuk menjaga ketahanan layanan saat
terjadi kegagalan simpul. Karakteristik \textit{durable log} dan pola
\textit{publish-subscribe} ini membuat Kafka sesuai sebagai
\textit{ingestion layer} pada Arsitektur Kappa, karena banyak konsumen dapat
mengolah aliran data yang sama secara independen \cite{kreps2014questioning}.
Dalam penelitian ini, Kafka berfungsi sebagai jalur distribusi hasil deteksi
dari modul visi komputer menuju pemrosesan \textit{streaming} secara kontinu,
sehingga analitik \textit{end-to-end} tetap berjalan konsisten
\cite{ayae2017video}.

\begin{figure}[H]
\centering
\includegraphics[width=0.8\textwidth]{gambar/kafka_arsitecture.png}
\caption{Arsitektur komponen inti Apache Kafka \cite{ayae2017video}}
\label{fig:kafka_architecture}
\end{figure}

\subsubsection{Apache Spark}
Apache Spark adalah \textit{unified analytics engine} untuk pemrosesan data
terdistribusi berskala besar pada mode \textit{batch} maupun \textit{stream}
\cite{zaharia2016apachespark}. Untuk kebutuhan \textit{real-time}, Spark
menyediakan {Spark Structured Streaming} yang memproses data dalam
bentuk \textit{incremental micro-batch} dengan semantik deklaratif
\cite{zaharia2018structured}. Pendekatan ini mendukung konsistensi hasil,
toleransi kesalahan melalui \textit{checkpointing}, serta
\textit{exactly-once processing semantics}. Selain itu, Spark Structured
Streaming mendukung pemrosesan \textit{stateful} untuk mempertahankan
informasi historis lintas aliran data. Mekanisme ini dapat diimplementasikan
melalui fungsi \textit{mapGroupsWithState} untuk membentuk \textit{state}
pada setiap grup data, serta \textit{flatMapGroupsWithState} untuk
memperbarui \textit{state} secara berkelanjutan.
% \subsubsection{Penghitungan Unik Berbasis Pelacakan}
% Dalam konteks pelacakan objek, terutama pada sistem CCTV, satu objek dapat muncul di banyak frame tanpa menunjukkan terjadinya pelanggaran baru. Oleh karena itu, penghitungan pelanggaran harus dilakukan berdasarkan \textit{identitas objek (ID tracking)}, bukan berdasarkan jumlah deteksi \cite{zhang2022bytetrack_eccv}.

% Dengan menggunakan \textit{stateful stream processing} pada Spark, sistem dapat menyimpan daftar \textit{ID} yang telah dihitung sehingga objek yang sama tidak dihitung lebih dari satu kali. Konsep ini dapat dimodelkan sebagai:

% \begin{equation}
% \mathbb{I}_{\text{count}}(ID) =
% \begin{cases}
% 1, & \text{jika } ID \notin \text{State}_{t-1} \\
% 0, & \text{jika } ID \in \text{State}_{t-1}
% \end{cases}
% \end{equation}

% Setelah ID dihitung sekali, state diperbarui:
% \begin{equation}
% \text{State}_{t} \leftarrow \text{State}_{t-1} \cup \{ID\}
% \end{equation}

% Dengan demikian, Spark Structured Streaming berperan sebagai modul pengelola state, memastikan bahwa penghitungan pelanggaran:
% \begin{enumerate}
%     \item Tidak mengalami \textit{double counting},
%     \item Dapat menangani objek yang terhalang sementara (occlusion),
%     \item Tetap konsisten meskipun data diterima secara kontinu dalam skala besar.
% \end{enumerate}

% \section{Object Detection}
% Object Detection adalah teknik dalam Computer Vision yang digunakan untuk mengidentifikasi dan menemukan objek tertentu dalam sebuah gambar atau video.
% Tujuan utama dari object detection adalah mengembangkan model dan teknik komputasi yang mampu mengidentifikasi jenis serta menentukan posisi objek dalam sebuah gambar.
% Secara umum, object detection dapat diklasifikasikan ke dalam dua bidang penelitian utama, yaitu general object detection dan detection application\cite{10028728}.
% General object detection berfokus pada pengembangan metode yang dapat mendeteksi berbagai objek dalam gambar secara menyeluruh, meniru cara kerja penglihatan manusia.
% Sementara itu, detection application lebih spesifik dalam mendeteksi objek dalam konteks tertentu, seperti identifikasi pejalan kaki, rambu lalu lintas, wajah, dan berbagai skenario lainnya \cite{10028728}.
\subsection{YOLO (\textit{You Only Look Once})}
\textit{You Only Look Once} (YOLO) adalah arsitektur deteksi objek berbasis
\textit{Convolutional Neural Network} (CNN) yang memprediksi
\textit{bounding box} dan kelas objek secara langsung dalam satu tahap
(\textit{single-stage detector}) \cite{redmon2016you}. Berbeda dengan metode
dua tahap seperti R-CNN, YOLO merumuskan deteksi sebagai masalah regresi
tunggal dari piksel citra ke koordinat kotak dan probabilitas kelas
\cite{redmon2016you}.

YOLO membagi citra masukan menjadi grid berukuran $S \times S$. Jika pusat
objek berada pada sebuah \textit{grid cell}, sel tersebut bertanggung jawab
melakukan prediksi objek. Setiap sel memprediksi $B$ \textit{bounding box}
dengan parameter $(x, y, w, h, \textit{confidence})$ serta probabilitas kelas.
Bentuk keluaran model dinyatakan pada Persamaan~\ref{eq:tensor_output}:
\begin{equation}
\label{eq:tensor_output}
\text{Tensor Output YOLO} \in \mathbb{R}^{S \times S \times (B \cdot 5 + C)}
\end{equation}
\noindent dengan:
\begin{description}
  \item[$\mathbb{R}$:] ruang bilangan real yang menyatakan dimensi tensor keluaran model;
  \item[$S \times S$:] jumlah sel \textit{grid} hasil pembagian citra masukan;
  \item[$B$:] jumlah \textit{bounding box} yang diprediksi tiap sel \textit{grid};
  \item[$5$:] jumlah parameter tiap \textit{bounding box}, yaitu $(x, y, w, h, \textit{confidence})$;
  \item[$C$:] jumlah kelas objek yang diprediksi tiap sel.
\end{description}
\begin{figure}[H]
\centering
\includegraphics[width=0.9\linewidth]{gambar/yolo_boundinbox.png}
\caption{Ilustrasi mekanisme prediksi \textit{bounding box} pada YOLO \cite{redmon2016you}}
\label{fig:yolo}
\end{figure}

Pada representasi ini, \(x, y\) menyatakan koordinat pusat kotak
yang diprediksi relatif terhadap sel grid tempat objek terdeteksi. Sementara
itu, \(w, h\) menyatakan lebar dan tinggi kotak pembatas relatif terhadap
ukuran citra masukan. Formulasi ini membantu menjaga konsistensi
lokalisasi ketika skala objek berubah dan mempermudah
konversi ke koordinat absolut pada tahap pascapemrosesan.

Skor \textit{confidence} didefinisikan sebagai kombinasi antara probabilitas
keberadaan objek dan ketepatan lokalisasi kotak pembatas. Nilai ini menjadi
indikator utama untuk menilai keyakinan prediksi pada setiap kandidat deteksi.
Semakin tinggi skor \textit{confidence}, semakin besar kemungkinan objek
benar-benar ada dan posisi kotaknya sesuai anotasi acuan. Pada tahap
pascapemrosesan, skor ini dipakai untuk menyaring prediksi bernilai rendah
sebelum proses \textit{non-maximum suppression} (NMS). Secara matematis,
\textit{confidence} dinyatakan pada Persamaan~\ref{eq:IOU_revised}:
\begin{equation}
C = \Pr(\text{Object}) \times \text{IoU}(\text{pred}, \text{truth})
\label{eq:IOU_revised}
\end{equation}
\noindent dengan:
\begin{description}
  \item[$C$:] skor \textit{confidence} sebuah kandidat deteksi;
  \item[$\Pr(\text{Object})$:] probabilitas keberadaan objek pada \textit{bounding box} yang diprediksi;
  \item[$\text{IoU}(\text{pred}, \text{truth})$:] nilai \textit{Intersection over Union} antara kotak prediksi (\textit{pred}) dan kotak anotasi acuan (\textit{truth}).
\end{description}

Fungsi kerugian YOLO menggunakan \textit{sum-squared error} untuk
menghitung kesalahan prediksi koordinat, ukuran kotak,
\textit{confidence}, dan klasifikasi \cite{redmon2016you}. Parameter
$\lambda_{\text{coord}}$ dan $\lambda_{\text{noobj}}$ digunakan untuk
mengatur bobot kesalahan koordinat dan penalti pada kotak tanpa objek. Bentuk
fungsi kerugian ditunjukkan pada Persamaan~\ref{eq:yolo_loss}:

\begin{equation}
\label{eq:yolo_loss}
\begin{split} \text{Loss} = & \lambda_{\text{coord}} \sum_{i=0}^{S^2} \sum_{j=0}^{B} \mathbb{I}_{ij}^{\text{obj}} \left[ (x_i - \hat{x}_i)^2 + (y_i - \hat{y}_i)^2 \right] \\ & + \lambda_{\text{coord}} \sum_{i=0}^{S^2} \sum_{j=0}^{B} \mathbb{I}_{ij}^{\text{obj}} \left[ (\sqrt{w_i} - \sqrt{\hat{w}_i})^2 + (\sqrt{h_i} - \sqrt{\hat{h}_i})^2 \right] \\ & + \sum_{i=0}^{S^2} \sum_{j=0}^{B} \mathbb{I}_{ij}^{\text{obj}} (C_i - \hat{C}_i)^2 \\ & + \lambda_{\text{noobj}} \sum_{i=0}^{S^2} \sum_{j=0}^{B} \mathbb{I}_{ij}^{\text{noobj}} (C_i - \hat{C}_i)^2 \\ & + \sum_{i=0}^{S^2} \mathbb{I}_{i}^{\text{obj}} \sum_{c \in \text{classes}} (p_i(c) - \hat{p}_i(c))^2
\end{split}
\end{equation}
\noindent dengan:
\begin{description}
  \item[$\lambda_{\text{coord}}$:] bobot galat untuk komponen koordinat dan ukuran kotak;
  \item[$\lambda_{\text{noobj}}$:] bobot penalti untuk sel yang tidak memuat objek;
  \item[$S^2$:] jumlah seluruh sel \textit{grid}, yaitu $S \times S$;
  \item[$B$:] jumlah \textit{bounding box} tiap sel \textit{grid};
  \item[$\mathbb{I}_{ij}^{\text{obj}}$:] indikator bernilai 1 jika \textit{box} ke-$j$ pada sel ke-$i$ bertanggung jawab atas objek, dan 0 jika sebaliknya;
  \item[$\mathbb{I}_{ij}^{\text{noobj}}$:] indikator bernilai 1 jika \textit{box} ke-$j$ pada sel ke-$i$ tidak memuat objek;
  \item[$\mathbb{I}_{i}^{\text{obj}}$:] indikator bernilai 1 jika objek berada pada sel ke-$i$;
  \item[$x_i, y_i, w_i, h_i$:] koordinat pusat, lebar, dan tinggi kotak acuan, sedangkan $\hat{x}_i, \hat{y}_i, \hat{w}_i, \hat{h}_i$ adalah nilai prediksinya;
  \item[$C_i$:] skor \textit{confidence} acuan pada sel ke-$i$, sedangkan $\hat{C}_i$ adalah nilai prediksinya;
  \item[$p_i(c)$:] probabilitas kelas $c$ acuan pada sel ke-$i$, sedangkan $\hat{p}_i(c)$ adalah nilai prediksinya;
  \item[$\text{classes}$:] himpunan seluruh kelas objek.
\end{description}

Arsitektur YOLO awal terdiri atas 24 lapisan konvolusional dan 2 lapisan
\textit{fully connected} \cite{redmon2016you}. Lapisan konvolusional
menggunakan \textit{filter} berukuran $3 \times 3$ dan $1 \times 1$ dengan
\textit{max pooling} untuk mereduksi dimensi spasial. Lapisan
\textit{fully connected} menghasilkan prediksi akhir kotak pembatas dan kelas
objek.
% \begin{figure}[H]
% \centering
% \includegraphics[width=0.8\textwidth]{gambar/yolo_architure.png}
% \caption{Arsitektur YOLO sebagai \textit{single-stage object detector}\cite{redmon2016you}.}
% \label{fig:yolov1_architecture}
% \end{figure}
% \subsubsection{Paradigma Prediksi Berbasis Anchor-Free}

% Pendekatan deteksi objek modern menunjukkan pergeseran konseptual dari metode \textit{anchor-based} menuju \textit{anchor-free prediction}.
% Pada varian terdahulu seperti YOLOv3 hingga YOLOv5, sistem deteksi bergantung pada kumpulan \textit{anchor boxes} yang telah ditentukan sebelumnya untuk memperkirakan ukuran dan rasio aspek objek.
% Meskipun efektif, mekanisme tersebut memiliki keterbatasan berupa ketergantungan pada hiperparameter dan sensitivitas terhadap variasi skala objek.

% YOLOv8 mengadopsi paradigma \textit{anchor-free} untuk meningkatkan efisiensi dan generalisasi model \cite{yaseen2024yolov8indepthexplorationinternal}.
% Dalam pendekatan ini, model tidak lagi melakukan regresi terhadap \textit{offset} dari \textit{anchor box} yang telah ditetapkan, tetapi secara langsung memprediksi parameter geometrik kotak pembatas yakni pusat, lebar, dan tinggi dalam koordinat normalisasi relatif terhadap peta fitur.
% Perubahan ini mengurangi kompleksitas komputasi, mempercepat konvergensi pelatihan, serta meningkatkan stabilitas pada kasus deteksi objek berukuran kecil dan padat.
% Strategi \textit{anchor-free} juga menunjukkan konsistensi hasil deteksi lintas domain tanpa memerlukan penyesuaian parameter \textit{anchor} secara manual.
% \begin{figure}[H]
%     \centering
%     \includegraphics[width=0.65\textwidth]{gambar/yolox-shifting-to-anchor-free-design.jpg}
%     \caption{Perbandingan pendekatan \textit{anchor-based} dan \textit{anchor-free}\cite{yaseen2024yolov8indepthexplorationinternal}.}
%     \label{fig:anchor_free_comparison}
% \end{figure}

% \subsubsection{Arsitektur Modul C2f}

% Modul C2f pada YOLOv8 merepresentasikan evolusi arsitektur dari modul \textbf{C3} yang digunakan pada YOLOv5, dengan tujuan meningkatkan efisiensi komputasi tanpa mengorbankan kemampuan representasi fitur.
% C2f mengimplementasikan prinsip \textit{Cross-Stage Partial (CSP)} yang memisahkan aliran fitur (\textit{feature flow}) menjadi dua jalur komputasi paralel: satu jalur dilewatkan melalui blok konvolusional bertingkat, sementara jalur lainnya langsung terhubung ke tahap penggabungan (\textit{feature concatenation}).
% Strategi ini mengurangi redundansi gradien selama propagasi balik, sehingga menurunkan jumlah parameter dan penggunaan memori secara signifikan \cite{yaseen2024yolov8indepthexplorationinternal}. Modul C3 yang hanya melakukan fusi fitur setelah proses konvolusi bertingkat, C2f memungkinkan \textit{feature fusion} terjadi secara lebih adaptif melalui koneksi parsial antar tahap, sehingga memperkaya konteks spasial yang ditangkap oleh jaringan.
% Pendekatan ini menghasilkan representasi fitur yang lebih stabil dan padat, sekaligus mempertahankan performa deteksi pada objek kecil dan berukuran menengah.
% Secara keseluruhan, studi perbandingan arsitektur menunjukkan bahwa C2f menawarkan peningkatan efisiensi sekitar 15–20\% dalam penggunaan memori dengan akurasi setara atau lebih baik dibandingkan modul sebelumnya \cite{selcuk2023comparison}.

% \begin{figure}[H]
%     \centering
%     \includegraphics[width=0.8\textwidth]{gambar/YOLOv5-C3-Module-YOLOv8-C2f-Module-and-Darknet-Bottleneck-Module.png}
%     \caption{Perbandingan struktur modul C3 (YOLOv5) dan C2f (YOLOv8). C2f meningkatkan efisiensi memori dengan mengoptimalkan pemrosesan jalur fitur sambil mempertahankan kedalaman representasi \cite{yaseen2024yolov8indepthexplorationinternal}.}
%     \label{fig:c3_vs_c2f}
% \end{figure}
% \subsubsection{Arsitektur \textit{Decoupled Head}}

% YOLOv8 memperkenalkan arsitektur \textit{Decoupled Head} sebagai perbaikan atas desain \textit{coupled head} yang digunakan pada varian sebelumnya.
% Pada pendekatan tradisional, jalur prediksi untuk \textit{kelas}, \textit{confidence}, dan \textit{bounding box regression} dibagikan dalam satu struktur konvolusional terpadu, yang sering kali menyebabkan interferensi gradien selama proses pelatihan.
% Fenomena ini menghambat konvergensi dan menurunkan stabilitas prediksi, khususnya pada skenario dengan kepadatan objek tinggi atau tumpang tindih antar objek.

% Arsitektur \textit{Decoupled Head} memisahkan jalur prediksi menjadi tiga cabang independen yang masing-masing difokuskan pada tugas klasifikasi, estimasi kepercayaan (\textit{objectness confidence}), dan regresi kotak pembatas.
% Setiap cabang melakukan pembelajaran terhadap representasi fitur yang terisolasi dan spesifik terhadap tujuannya \cite{yaseen2024yolov8indepthexplorationinternal}.
% Pemisahan ini menghasilkan distribusi gradien yang lebih seimbang antar cabang prediksi, yang berdampak pada peningkatan stabilitas pelatihan dan percepatan konvergensi.
% Hasil empiris menunjukkan bahwa struktur tersebut juga memperkuat keseimbangan antara presisi dan \textit{recall} pada evaluasi \textit{mAP}, terutama pada dataset dengan tingkat tumpang tindih objek yang tinggi.

% \begin{figure}[H]
%     \centering
%     \includegraphics[width=0.65\textwidth]{gambar/The-difference-between-the-YOLO-coupled-head-and-decoupled-head.png}
%     \caption{Perbandingan arsitektur \textit{coupled head} dan \textit{decoupled head} pada YOLOv8.
%     Jalur regresi kotak, klasifikasi, dan estimasi kepercayaan dipisahkan untuk menghindari interferensi gradien, meningkatkan stabilitas pelatihan, serta mempercepat konvergensi model \cite{yaseen2024yolov8indepthexplorationinternal}.}
%     \label{fig:decoupled_head_yolov8}
% \end{figure}

\subsection{ByteTrack: Multi-Object Tracking Berbasis Asosiasi Dua Tahap}

Pada sistem video \textit{real-time}, identitas objek perlu dipertahankan
secara konsisten antar-frame agar pelacakan berkelanjutan. Deteksi objek saja
umumnya hanya menghasilkan posisi objek pada satu frame tanpa menjaga
identitas lintas waktu. Kebutuhan ini menjadi alasan penggunaan pendekatan
\textit{Multi-Object Tracking} (MOT), salah satunya \textit{ByteTrack}
\cite{zhang2022bytetrack_eccv}, yang menerapkan asosiasi dua tahap untuk
menjaga kontinuitas identitas objek saat bergerak. ByteTrack mempertimbangkan deteksi berskor rendah yang kerap muncul akibat
\textit{occlusion}, perubahan pencahayaan, atau ukuran objek kecil. Asosiasi
dilakukan secara bertahap, pertama pada deteksi skor tinggi, lalu pada deteksi
skor rendah untuk memulihkan identitas yang sempat hilang. Strategi ini
membantu menjaga pelacakan dan mengurangi risiko kehilangan identitas
pada lingkungan dinamis \cite{zhang2022bytetrack_eccv}.

\begin{figure}[H]
\centering
\includegraphics[width=0.6\textwidth]{gambar/bytraking_pipeline.png}
\caption{Pipeline ByteTrack \cite{zhang2022bytetrack_eccv}. Arsitektur ini mengintegrasikan tahap deteksi dan asosiasi secara bertingkat untuk mempertahankan identitas objek pada video \textit{real-time}.}
\label{fig:bytetrack_pipeline}
\end{figure}

\subsubsection{Asosiasi Dua Tahap}

ByteTrack menerapkan strategi asosiasi dua tahap untuk mempertahankan
identitas objek secara konsisten antar-frame. Hasil deteksi pada setiap frame
dikelompokkan ke dalam dua kategori, yaitu
deteksi skor tinggi dan deteksi skor rendah. Prediksi posisi objek diperoleh
dengan \textit{Kalman Filter}, sedangkan pencocokan antara deteksi dan prediksi
dilakukan menggunakan metrik \textit{Intersection over Union} (IoU).

Pada tahap pertama, lintasan aktif dicocokkan dengan deteksi skor tinggi untuk
mendapatkan pasangan yang paling sesuai. Tahap kedua meninjau lintasan yang belum
terasosiasi dan mencocokkannya dengan deteksi skor rendah agar identitas dapat
dipulihkan. Penugasan akhir diselesaikan menggunakan
\textit{Hungarian Algorithm} untuk meminimalkan biaya asosiasi berbasis IoU.
Pendekatan ini membantu meningkatkan stabilitas pelacakan pada kondisi
\textit{occlusion} dan variasi pencahayaan tinggi.

\begin{figure}[H]
\centering
\includegraphics[width=0.8\textwidth]{gambar/dua_langkah.png}
\caption{Ilustrasi proses asosiasi dua tahap pada ByteTrack. Tahap pertama mengasosiasikan deteksi skor tinggi, kemudian tahap kedua memulihkan identitas objek melalui asosiasi dengan deteksi skor rendah \cite{zhang2022bytetrack_eccv}.}
\label{fig:bytetrack_two_stage}
\end{figure}

\subsection{Asosiasi Spasial Berbasis \textit{Intersection over Area} (IoA)}
\label{subsec:ioa}

Deteksi pelanggaran helm tidak cukup hanya mengenali pengendara tanpa helm,
tetapi juga harus memastikan objek tersebut berada pada sepeda motor. Kebutuhan
ini menuntut mekanisme asosiasi spasial antara \textit{bounding box} pengendara
dan \textit{bounding box} sepeda motor \cite{aboah2023helmet, shoman2024enforcing}.
Metrik tumpang tindih kotak yang umum dipakai adalah \textit{Intersection over
Union} (IoU), yaitu rasio luas irisan terhadap luas gabungan dua kotak
\cite{lin2014microsoft}. Akan tetapi, IoU kurang sesuai ketika kedua kotak
berbeda ukuran secara mencolok, seperti kotak pengendara yang jauh lebih kecil
daripada kotak sepeda motor. Pada kondisi tersebut, nilai IoU menjadi rendah
meskipun pengendara berada di dalam area sepeda motor.

\textit{Intersection over Area} (IoA) adalah varian IoU yang menormalkan luas
irisan terhadap luas salah satu kotak, bukan terhadap luas gabungan. Pada
asosiasi pengendara--motor, luas irisan dibagi dengan luas kotak pengendara
sebagai kotak acuan, sehingga nilai IoA menyatakan proporsi area kotak
pengendara yang berada di dalam kotak sepeda motor. Bentuk umum IoA antara kotak
acuan $A$ dan kotak pembanding $B$ dinyatakan pada Persamaan~\ref{eq:ioa-umum}:
\begin{equation}
\label{eq:ioa-umum}
\text{IoA}(A,B) = \frac{|A \cap B|}{|A|}
\end{equation}
\noindent dengan:
\begin{description}
  \item[$A$:] kotak pembatas acuan yang luasnya menjadi penyebut, yaitu kotak objek berukuran lebih kecil (pengendara);
  \item[$B$:] kotak pembatas pembanding, yaitu kotak objek berukuran lebih besar (sepeda motor);
  \item[$|A \cap B|$:] luas area irisan antara kotak $A$ dan kotak $B$, dalam satuan piksel;
  \item[$|A|$:] luas area kotak acuan $A$, dalam satuan piksel.
\end{description}

Nilai IoA berada pada rentang $[0,1]$. Nilai mendekati 1 menunjukkan kotak acuan
hampir sepenuhnya berada di dalam kotak pembanding, sedangkan nilai mendekati 0
menunjukkan kedua kotak hampir tidak beririsan. Berbeda dengan IoU yang nilainya
menyusut ketika selisih ukuran kotak besar, IoA tetap tinggi selama kotak kecil
hampir seluruhnya berada di dalam kotak besar. Sifat ini membuat IoA lebih tepat
untuk memverifikasi keterkaitan pengendara dengan sepeda motor pada kamera
statis bersudut pandang menyamping \cite{shoman2024enforcing}. Ambang IoA
ditetapkan secara empiris sebagai batas keputusan asosiasi; rinciannya dibahas
pada Bab~\ref{chap:metodologi}.

\subsection{Data Warehouse}

\textit{Data warehouse} adalah sistem penyimpanan terpusat yang dirancang
untuk mendukung analitik dan pengambilan keputusan jangka panjang. Berbeda
dengan basis data operasional yang berfokus pada transaksi harian,
\textit{data warehouse} dioptimalkan untuk \textit{query} kompleks, agregasi,
dan analisis historis \cite{vyas2025demystifying}. Struktur ini memungkinkan
integrasi data dari berbagai sumber sehingga informasi strategis dapat
dihasilkan secara lebih lengkap. Data di dalam \textit{data warehouse}
bersifat terintegrasi, berorientasi subjek, historis, dan relatif stabil.
Informasi disusun untuk mendukung analisis multidimensi, pembacaan tren
jangka panjang, serta perbandingan antarkurun waktu. Dalam praktik saat ini,
konsep ini juga digunakan pada
lingkungan \textit{cloud} dan \textit{big data} untuk meningkatkan
skalabilitas serta efisiensi analitik \cite{golfarelli2024bigdata}.
\subsection{Data Mart}

\textit{Data mart} merupakan subset dari \textit{data warehouse} yang
dirancang untuk kebutuhan analitik pada domain spesifik, bukan untuk
pemrosesan transaksi harian (\textit{online transaction processing}/OLTP).
Dalam konteks sistem pemantauan pelanggaran helm, rancangan ini memungkinkan
penetapan \textit{event grain} yang jelas pada tabel fakta. Dengan demikian,
setiap kejadian dapat terekam konsisten dan tetap dapat diaudit. Penggunaan
tabel dimensi yang dapat digunakan ulang, seperti dimensi kamera, jenis
pelanggaran, dan waktu, juga menyediakan struktur deskriptif untuk
analisis dari beberapa sudut pandang.

Dari sisi performa, \textit{data mart} menyediakan jalur agregasi yang
lebih efisien untuk kebutuhan panel Grafana karena metrik telah dipetakan pada
skema relasional yang terarah. Struktur ini juga berfungsi sebagai lapisan
definisi yang konsisten untuk definisi indikator kinerja utama (KPI), sehingga
istilah, rumus, dan konteks analitik tetap konsisten pada seluruh dashboard.
Tanpa lapisan \textit{data mart}, kueri analitik harus berulang kali
merekonstruksi logika bisnis langsung dari aliran data mentah. Kondisi ini
dapat meningkatkan latensi, menghasilkan inkonsistensi, dan menambah kompleksitas
pemeliharaan sistem.

\subsubsection{Star Schema}

\textit{Star schema} merupakan pendekatan pemodelan data multidimensi yang
banyak diterapkan pada sistem \textit{data warehouse}. Struktur ini terdiri
atas satu tabel fakta yang menyimpan metrik numerik serta tabel dimensi yang
menyediakan konteks deskriptif. Desain ini mendukung efisiensi analitik
melalui denormalisasi terkontrol, sehingga \textit{query} agregatif dan
pelaporan dapat dijalankan lebih cepat. Kombinasi tabel fakta dan dimensi juga
memberi fleksibilitas untuk analisis multi-sumbu serta pelaporan berbasis
waktu.

Pendekatan \textit{star schema} memiliki beberapa keunggulan:
\begin{enumerate}[a.]
\item Memungkinkan eksekusi \textit{query} analitik berkecepatan tinggi melalui struktur denormalisasi,
\item Menyediakan skema data yang mudah diinterpretasikan untuk pelaporan dan visualisasi,
\item Mendukung analisis \textit{time-series} serta perbandingan historis secara efisien.
\end{enumerate}

\subsection{Grafana}

\textit{Grafana} adalah platform sumber terbuka untuk visualisasi data dan
pengembangan \textit{dashboard} interaktif pada lingkungan
\textit{big data warehouse}. Platform ini dapat terhubung dengan berbagai sistem,
seperti \textit{Apache Kafka}, \textit{Prometheus}, \textit{InfluxDB}, dan
\textit{Elasticsearch}, sehingga sesuai untuk pemantauan \textit{real-time}
\cite{sutanto2025wazuhgrafana}. Melalui antarmuka web, Grafana membantu pengguna
memantau performa sistem, mendeteksi anomali, dan melakukan analisis
operasional .

Grafana berperan sebagai lapisan presentasi yang menghubungkan hasil
pemrosesan sistem analitik, seperti \textit{Apache Spark}, ke dalam
visualisasi interaktif yang diperbarui otomatis. Aliran data dari
\textit{streaming pipeline} diproses melalui mekanisme
\textit{Extract--Transform--Load} (ETL) sebelum ditampilkan pada dashboard.
Rancangan ini mendukung \textit{real-time analytics}, sehingga
pengguna dapat mengeksplorasi data terbaru tanpa menunggu siklus
\textit{batch} \cite{devireddy2020realtimewarehouse}.

\subsection{Evaluasi Deteksi Objek}

Evaluasi model deteksi dilakukan untuk menilai kemampuan model mengenali dan
melokalisasi objek dibandingkan anotasi \textit{ground truth}. Proses ini
mengikuti protokol \textit{MS COCO} \cite{lin2014microsoft} dengan metrik
utama \textit{Precision}, \textit{Recall}, \textit{Average Precision} (AP),
dan \textit{Mean Average Precision} (mAP). Melalui metrik tersebut,
keseimbangan antara deteksi benar dan kesalahan prediksi dapat dinilai secara
menyeluruh.

\subsubsection{Klasifikasi Deteksi: TP, FP, FN, dan TN}

Formulasi berbasis \textit{confusion matrix} menggunakan elemen $TC_{ij}$
untuk menyatakan jumlah sampel dengan kelas sebenarnya $i$ yang diprediksi
sebagai kelas $j$. Pada kasus multikelas dengan $K$ kelas, nilai diagonal
$TC_{ii}$ menunjukkan prediksi benar untuk kelas $i$, sedangkan elemen di luar
diagonal menunjukkan kesalahan antar kelas. Komponen evaluasi untuk kelas
target $i$ dituliskan sebagai:
\begin{equation}
\begin{aligned}
TP_i &= TC_{ii}, & FP_i &= \sum_{j \neq i} TC_{ji}, \\
FN_i &= \sum_{j \neq i} TC_{ij}, & TN_i &= \sum_{p \neq i} \sum_{q \neq i} TC_{pq}
\end{aligned}
\end{equation}
\noindent dengan:
\begin{description}
  \item[$TC_{ij}$:] jumlah sampel berkelas sebenarnya $i$ yang diprediksi sebagai kelas $j$;
  \item[$K$:] jumlah kelas pada kasus multikelas;
  \item[$i$:] indeks kelas target yang sedang dianalisis;
  \item[$j$:] indeks kelas selain kelas target ($j \neq i$);
  \item[$p, q$:] indeks penjumlahan baris dan kolom selain kelas target pada perhitungan $TN_i$;
  \item[$TP_i, FP_i, FN_i, TN_i$:] berturut-turut jumlah \textit{true positive}, \textit{false positive}, \textit{false negative}, dan \textit{true negative} untuk kelas $i$.
\end{description}

Makna setiap metrik adalah sebagai berikut. \textit{True Positive} (TP) adalah
jumlah prediksi benar pada kelas yang sedang dianalisis. \textit{False
Positive} (FP) adalah sampel dari kelas lain yang salah diprediksi sebagai
kelas target. \textit{False Negative} (FN) adalah sampel kelas target yang
salah diprediksi ke kelas lain. \textit{True Negative} (TN) adalah seluruh
sampel non-target yang juga tidak diprediksi sebagai kelas target. Pada
\textit{object detection}, perhitungan ini dilakukan setelah pencocokan kelas
dan ambang IoU. Deteksi dihitung TP hanya jika kelas tepat dan lokalisasi
memenuhi ambang.

Pada skenario klasifikasi biner (kelas positif dan negatif), komponen evaluasi
didefinisikan sebagai berikut: TP merepresentasikan prediksi positif yang benar
(positif$\to$positif), FP merepresentasikan prediksi positif yang keliru pada
sampel negatif (negatif$\to$positif), FN merepresentasikan sampel positif yang
salah diprediksi sebagai negatif (positif$\to$negatif), dan TN
merepresentasikan prediksi negatif yang benar (negatif$\to$negatif). Pada
skenario klasifikasi tiga kelas (A, B, C) dengan kelas A sebagai target,
komponen evaluasi dinyatakan sebagai TP$_A$ = $TC_{AA}$,
FP$_A$ = $TC_{BA}+TC_{CA}$, FN$_A$ = $TC_{AB}+TC_{AC}$, sedangkan TN$_A$
mencakup seluruh elemen yang tidak berada pada baris A maupun kolom A.
Formulasi tersebut menjadi landasan perhitungan \textit{Precision} dan
\textit{Recall} per kelas sebelum dilakukan agregasi ke metrik evaluasi
keseluruhan.

\begin{figure}[H]
\centering
\includegraphics[width=0.9\textwidth]{gambar/false_positive.png}
\caption{Ilustrasi kategori hasil deteksi objek pada sistem \textit{object detection}: True Positive (TP), False Positive (FP), dan False Negative (FN) \cite{lin2014microsoft}.}
\label{fig:tp_fp_fn_tn}
\end{figure}

Klasifikasi ini menjadi dasar kuantitatif untuk menilai keseimbangan
antara kemampuan model mendeteksi objek relevan dan risiko kesalahan deteksi.
Relasi antarkomponen ini digunakan untuk menafsirkan metrik lanjutan,
seperti kurva \textit{Precision-Recall} dan mAP. Relasi antara
\textit{Precision}, \textit{Recall}, dan \textit{Intersection over Union}
(IoU) ditunjukkan pada Gambar~\ref{fig:precision_recall_iou}.

Perhitungan \textit{Precision} dan \textit{Recall} dinyatakan sebagai:
\begin{equation}
\text{Precision} = \frac{TP}{TP + FP}, \qquad
\text{Recall} = \frac{TP}{TP + FN}
\end{equation}
\noindent dengan:
\begin{description}
  \item[$TP$:] jumlah \textit{true positive}, yaitu deteksi positif yang benar;
  \item[$FP$:] jumlah \textit{false positive}, yaitu deteksi positif yang keliru;
  \item[$FN$:] jumlah \textit{false negative}, yaitu objek target yang tidak terdeteksi.
\end{description}

\textit{Intersection over Union} (IoU) digunakan untuk mengukur tingkat
kesesuaian geometrik antara kotak prediksi dan anotasi
\textit{ground truth}, yang dirumuskan sebagai:
\begin{equation}
\text{IoU}(A,B) = \frac{|A \cap B|}{|A \cup B|}
\end{equation}
\noindent dengan:
\begin{description}
  \item[$A$:] kotak pembatas prediksi;
  \item[$B$:] kotak pembatas anotasi acuan (\textit{ground truth});
  \item[$|A \cap B|$:] luas area irisan antara kotak $A$ dan kotak $B$;
  \item[$|A \cup B|$:] luas area gabungan kotak $A$ dan kotak $B$.
\end{description}

Nilai IoU yang lebih tinggi menunjukkan tumpang tindih area yang lebih besar
antara prediksi dan anotasi, sehingga ketepatan lokalisasi objek menjadi lebih
baik. Kombinasi evaluasi berbasis \textit{Precision}, \textit{Recall}, dan IoU
kemudian digunakan untuk menghitung metrik agregat, terutama AP dan mAP.

\begin{figure}[H]
\centering
\includegraphics[width=0.9\textwidth]{gambar/object-detection-metrics_3.png}
\caption{Hubungan antara \textit{Precision}, \textit{Recall}, dan \textit{Intersection over Union} (IoU) pada sistem deteksi objek. Nilai IoU dihitung dari area tumpang tindih antara kotak prediksi dan anotasi \textit{ground truth} sesuai protokol evaluasi deteksi objek \cite{lin2014microsoft}.}
\label{fig:precision_recall_iou}
\end{figure}

\subsubsection{Average Precision (AP) dan Mean Average Precision (mAP)}

Evaluasi kinerja deteksi objek umumnya menggunakan metrik
\textit{Average Precision} (AP) dan \textit{Mean Average Precision} (mAP),
yang menjadi standar utama pada dataset \textit{MS COCO}
\cite{lin2014microsoft}. Kedua metrik ini menilai keseimbangan antara
ketepatan prediksi kelas dan ketelitian lokalisasi terhadap
\textit{ground truth}.

\textit{Average Precision} (AP) didefinisikan sebagai luas area di bawah kurva
\textit{Precision-Recall} (PR) untuk setiap kelas objek. Nilai AP
menggambarkan rata-rata presisi pada berbagai tingkat \textit{recall}, dan
secara matematis dinyatakan sebagai:
\begin{equation}
\text{AP} = \int_{0}^{1} \text{Precision}(\text{Recall}) \, d(\text{Recall})
\end{equation}
\noindent dengan:
\begin{description}
  \item[$\text{AP}$:] \textit{Average Precision}, yaitu luas area di bawah kurva \textit{Precision-Recall} satu kelas;
  \item[$\text{Precision}(\text{Recall})$:] nilai presisi sebagai fungsi dari \textit{recall};
  \item[$d(\text{Recall})$:] elemen integrasi terhadap \textit{recall} pada rentang $[0,1]$.
\end{description}
Nilai AP yang tinggi berarti model mampu mempertahankan presisi
tinggi pada rentang \textit{recall} yang luas. Untuk menilai performa model
secara keseluruhan lintas kelas dan ambang lokalisasi, digunakan mAP yang
merupakan rata-rata AP pada beberapa ambang IoU. Formulasi mAP pada
\textit{COCO benchmark} adalah:
\begin{equation}
\text{mAP}_{@[0.50:0.95]} = \frac{1}{10} \sum_{k=0}^{9} \text{AP}_{\text{IoU} = 0.50 + 0.05k}
\end{equation}
\noindent dengan:
\begin{description}
  \item[$\text{mAP}_{@[0.50:0.95]}$:] rata-rata AP pada sepuluh ambang IoU dari 0{,}50 hingga 0{,}95;
  \item[$k$:] indeks penjumlahan ambang IoU, bernilai 0 sampai 9;
  \item[$\text{AP}_{\text{IoU} = 0.50 + 0.05k}$:] nilai AP pada ambang IoU sebesar $0{,}50 + 0{,}05k$;
  \item[$\tfrac{1}{10}$:] faktor perata-rataan atas sepuluh ambang IoU.
\end{description}

Pendekatan ini mencakup sepuluh ambang IoU dari 0{.}50 hingga 0{.}95 dengan
interval 0{.}05, sehingga evaluasi tidak hanya menilai keberhasilan deteksi
objek, tetapi juga ketepatan batas spasialnya.
% \subsubsection{MOTA (Multi-Object Tracking Accuracy)}
% MOTA merupakan salah satu metrik paling luas digunakan untuk mengevaluasi kinerja pelacakan karena menggambarkan akurasi pelacakan secara keseluruhan.
% MOTA didefinisikan sebagai:
% \begin{equation}
% \text{MOTA} = 1 - \frac{\sum_{t} (FN_t + FP_t + IDSW_t)}{\sum_{t} GT_t}
% \end{equation}
% \noindent
% dengan:
% \begin{description}
%     \item[$FN_t$:] jumlah objek yang tidak terdeteksi pada frame ke-$t$,
%     \item[$FP_t$:] jumlah deteksi palsu,
%     \item[$IDSW_t$:] jumlah \textit{ID Switch} (pergantian identitas),
%     \item[$GT_t$:] jumlah total objek sebenarnya (ground truth).
% \end{description}
% \begin{figure}[H]
%     \centering
%     \includegraphics[width=0.75\textwidth]{gambar/mota.png}
%     \caption{Ilustrasi MOTA.
%     Gambar menunjukkan hubungan antara deteksi benar, salah, serta pergantian identitas dalam penilaian MOTA.}
%     \label{fig:mota_milan}
% \end{figure}
\subsection{Evaluasi Sistem Streaming}

Evaluasi sistem pemrosesan data \textit{real-time} berbasis arsitektur Kappa
tidak hanya menilai akurasi keluaran, tetapi juga efisiensi aliran data yang
berjalan kontinu dengan jeda rendah
\cite{kreps2014questioning}. Dua metrik utama yang digunakan adalah
\textit{throughput} dan \textit{latency}, yang bersama-sama menggambarkan
kinerja \textit{pipeline} pemrosesan \textit{end-to-end}
\cite{kleppmann2017designing}. \textit{Throughput} menunjukkan jumlah unit
data, seperti \textit{frame} video atau pesan Kafka, yang dapat diproses
sistem per satuan waktu. Metrik ini umumnya diukur dalam
\textit{frame per second} (FPS) \cite{marz2015big}.

\textit{latency} menggambarkan jeda antara data diterima dan hasil
dikeluarkan \cite{kleppmann2017designing}. Pada arsitektur Kappa, seluruh
pemrosesan berlangsung dalam satu jalur \textit{stream pipeline}. Karena itu,
nilai \textit{latency} dipengaruhi oleh rancangan pemrosesan
\textit{stateful} dan efisiensi komunikasi antarkomponen. Sistem \textit{streaming}
dinilai baik jika mampu mempertahankan \textit{throughput} tinggi dengan
\textit{latency} rendah, idealnya pada rentang \textit{sub-second}.
